\documentclass[10pt, a4paper]{article}

\usepackage[utf8]{inputenc}
\usepackage[T1]{fontenc}
\usepackage[ngerman]{babel}
\usepackage[margin=1in]{geometry}
\usepackage{xcolor}
\usepackage{hyperref}
\usepackage{charter}

% Define Primary Color: RGB(0, 51, 102)
\definecolor{primarycolor}{RGB}{0, 51, 102}

\hypersetup{
    colorlinks=true,
    linkcolor=primarycolor,
    urlcolor=primarycolor,
}

\pagestyle{empty}

\begin{document}

% SENDER INFO
\noindent
\textbf{\large Kartavya Niraj Dikshit} \\
Halbmondstraße 7 \\
91054 Erlangen \\
Mobil: +49 155 10940829 \\
E-Mail: \href{mailto:kartavya.dikshit@gmail.com}{kartavya.dikshit@gmail.com} \\
LinkedIn: \href{https://www.linkedin.com/in/kartavya-dikshit}{linkedin.com/in/kartavya-dikshit} \\
GitHub: \href{https://github.com/KartavyaDikshit}{github.com/KartavyaDikshit}

\vspace{1cm}

% RECIPIENT INFO
\noindent
Siemens Healthineers AG \\
Talent Acquisition Team \\
Henkestraße 127 \\
91052 Erlangen

\vspace{1cm}

% DATE AND SUBJECT
\hfill Erlangen, den 19. März 2026

\vspace{0.5cm}

\noindent
\textbf{\large Bewerbung als Werkstudent (f/m/d) Full-Stack Software Development for Innovation Projects} \\
\textbf{Referenznummer: 498493}

\vspace{1cm}

% SALUTATION
\noindent
Sehr geehrte Damen und Herren,

\vspace{0.5cm}

% CONTENT
\noindent
als Masterstudent der Data Science an der Friedrich-Alexander-Universität Erlangen-Nürnberg mit einer tiefen Spezialisierung auf Künstliche Intelligenz und Machine Learning verfolge ich die Innovationen von Siemens Healthineers im Bereich der digitalen Gesundheitsversorgung mit großer Begeisterung. Ihr Pioniergeist, modernste Technologie mit medizinischer Exzellenz zu verbinden, motiviert mich, mein technisches Know-how in Ihr Team für Innovationsprojekte einzubringen.

In meinem bisherigen Werdegang konnte ich fundierte Erfahrungen in der Full-Stack-Entwicklung sammeln, die exakt Ihr Anforderungsprofil widerspiegeln. In meinen Projekten habe ich skalierbare App-Prototypen mittels \textbf{React Native (Expo)} sowie robuste Backend-Lösungen mit \textbf{Node.js und Python} realisiert. Besonders hervorheben möchte ich meine Fähigkeit, komplexe Stakeholder-Anforderungen in präzise technische Arbeitspakete zu übersetzen – eine Kompetenz, die ich in agilen Entwicklungsumgebungen (Scrum/Kanban) unter Einsatz von \textbf{Git, Docker und Azure DevOps} erfolgreich unter Beweis gestellt habe.

Mein aktuelles Studium an der FAU ermöglicht es mir, neueste wissenschaftliche Erkenntnisse aus dem Bereich der KI-gestützten Datenanalyse direkt in die Praxis zu transferieren. Die Herausforderung, innovative App-Prototypen für den Gesundheitssektor zu entwickeln, reizt mich besonders, da hier technologische Präzision und gesellschaftlicher Mehrwert unmittelbar aufeinandertreffen. Durch meinen Wohnort in Erlangen bin ich zudem zeitlich flexibel und kann die Zusammenarbeit an den Standorten Erlangen und Forchheim problemlos realisieren.

Ich bin davon überzeugt, dass ich durch meine strukturierte Arbeitsweise, meine schnelle Auffassungsgabe für neue Technologien und meine Leidenschaft für Full-Stack-Lösungen einen wertvollen Beitrag zu Ihren Innovationsprojekten leisten werde.

Gerne überzeuge ich Sie in einem persönlichen Gespräch von meiner Eignung. Ich stehe Ihnen ab sofort für mindestens ein Jahr mit einer wöchentlichen Arbeitszeit von 15 Stunden zur Verfügung.

\vspace{1cm}

% SIGN-OFF
\noindent
Mit freundlichen Grüßen

\vspace{1.5cm}

\noindent
Kartavya Niraj Dikshit
\end{document}